%% Shared body of the math-accent specimen.
%%
%% Not a document root -- it is \input by gwbook-accents.tex and
%% gwbook-accents-charter.tex, which supply the \documentclass and differ only in
%% the typeface. It lives in include/ because `make latex` treats every
%% top-level .tex in latex/ as a root and would otherwise compile this alone.
%%
%% This sheet exists to be looked at: the class re-centres math accents over
%% every upright glyph (see the "Math accents" note in gwbook.cls), and the only
%% way to know that worked for all 52 letters in both faces is to print them.

%% ---------------------------------------------------------------------------
%% The hairline used in "Reading the rule" below.
%% ---------------------------------------------------------------------------
%% The rule is drawn at the centre of the base letter's ink -- which is a glyph
%% property, not a box property, so TeX cannot measure it and Lua has to. Walk
%% the typeset box, find the base glyph (the first one that is not a combining
%% accent), and report the middle of its bounding box.

\newdimen\gwbcentre

\begin{luacode*}
local GLYPH, HLIST, VLIST = node.id("glyph"), node.id("hlist"), node.id("vlist")

local function is_accent(cp)
  return (cp >= 0x0300 and cp <= 0x036F) or (cp >= 0x20D0 and cp <= 0x20EF)
end

-- Returns the base glyph and its x, in sp from the left edge of the list.
-- Inside a vlist the children all start at the same x and a shift moves them
-- right; inside an hlist each one moves x along by its own width.
local function base_glyph(head, x, vertical)
  for n in node.traverse(head) do
    if n.id == GLYPH then
      if not is_accent(n.char) then return n, x end
      if not vertical then x = x + n.width end
    elseif n.id == HLIST or n.id == VLIST then
      local glyph, gx = base_glyph(n.head, x + (vertical and n.shift or 0),
                                   n.id == VLIST)
      if glyph then return glyph, gx end
      if not vertical then x = x + n.width end
    elseif not vertical and n.width then
      x = x + n.width
    end
  end
end

function gwb_ink_centre(boxnum)
  local centre = 0
  local glyph, x = base_glyph(tex.box[boxnum].head, 0, false)
  if glyph then
    local tfm = fonts.hashes.identifiers[glyph.font]
    local desc = tfm.descriptions and tfm.descriptions[glyph.char]
    local bbox = desc and desc.boundingbox
    if bbox then
      centre = x + (bbox[1] + bbox[3]) / 2 * tfm.size / (tfm.units or 1000)
    else
      centre = x + glyph.width / 2
    end
  end
  tex.setdimen("gwbcentre", centre)
end
\end{luacode*}

\definecolor{gwbMark}{HTML}{C03030}

% \ruled{<math>} -- the sample with a hairline through the base letter's ink
% centre. \rlap has no width, so the rule and the sample share an origin.
\newcommand{\ruled}[1]{%
  \setbox0=\hbox{$#1$}%
  \directlua{gwb_ink_centre(0)}%
  \leavevmode
  \rlap{\kern\gwbcentre
        \textcolor{gwbMark}{\vrule width 0.2pt height 3.1ex depth 0.9ex}}%
  \box0}

% \gwbcaps{\symbf}{label} prints \hat over every capital in that alphabet,
% \gwblows the same over the lowercase; the starred forms add the hairline.
% \foreach comes from pgffor, which tikz brings in.
\newlength{\gwblabelwidth}\setlength{\gwblabelwidth}{4.6em}
\newcommand{\gwblabel}[1]{\makebox[\gwblabelwidth][l]{\footnotesize\ttfamily#1}}
% A row that runs past the measure wraps under its first letter rather than
% back to the margin, so the alphabet still reads as one block.
\newcommand{\gwbrowstart}[1]{%
  \par\nobreak\vspace{2pt}\noindent
  \hangindent\gwblabelwidth \hangafter 1 \gwblabel{#1}}
\newcommand{\gwbrow}[3]{%
  \gwbrowstart{#2}%
  \foreach \gwbL in {#3}{$\hat{#1{\gwbL}}$\hskip 0.3em}\par}
\newcommand{\gwbrowruled}[3]{%
  \gwbrowstart{#2}%
  \foreach \gwbL in {#3}{\ruled{\hat{#1{\gwbL}}}\hskip 0.3em}\par}
% The Greek rows are spelled out rather than looped: \foreach wants a list of
% characters, and these are control sequences.
\newcommand{\gwbg}[2]{$\hat{#1{#2}}$\hskip 0.3em}
\newcommand{\gwbgreek}[1]{%
  \gwbg{#1}{\Gamma}\gwbg{#1}{\Delta}\gwbg{#1}{\Theta}\gwbg{#1}{\Lambda}%
  \gwbg{#1}{\Xi}\gwbg{#1}{\Pi}\gwbg{#1}{\Sigma}\gwbg{#1}{\Upsilon}%
  \gwbg{#1}{\Phi}\gwbg{#1}{\Psi}\gwbg{#1}{\Omega}\hskip 0.6em
  \gwbg{#1}{\alpha}\gwbg{#1}{\beta}\gwbg{#1}{\gamma}\gwbg{#1}{\delta}%
  \gwbg{#1}{\epsilon}\gwbg{#1}{\zeta}\gwbg{#1}{\eta}\gwbg{#1}{\theta}%
  \gwbg{#1}{\iota}\gwbg{#1}{\kappa}\gwbg{#1}{\lambda}\gwbg{#1}{\mu}%
  \gwbg{#1}{\nu}\gwbg{#1}{\xi}\gwbg{#1}{\pi}\gwbg{#1}{\rho}%
  \gwbg{#1}{\sigma}\gwbg{#1}{\tau}\gwbg{#1}{\upsilon}\gwbg{#1}{\phi}%
  \gwbg{#1}{\chi}\gwbg{#1}{\psi}\gwbg{#1}{\omega}\par}

\newcommand{\gwbcaps}[2]{\gwbrow{#1}{#2}{A,...,Z}}
\newcommand{\gwblows}[2]{\gwbrow{#1}{#2}{a,...,z}}
\newcommand{\gwbcapsruled}[2]{\gwbrowruled{#1}{#2}{A,...,Z}}
\newcommand{\gwblowsruled}[2]{\gwbrowruled{#1}{#2}{a,...,z}}

\maketitle
\tableofcontents

\chapter{Math accents}
\label{ch:accents}

An accent is positioned by a single number the font stores for each glyph: the
top accent attachment point, the x where the peak of the accent goes. Font
designers set it with care for the slanted alphabets, where an accent belongs
over the top of the letter rather than over the middle of its ink. For the
upright alphabets they often do not --- Schola's bold \(\symbf{R}\) carried a
value about \(1.2\,\text{pt}\) left of centre, and XCharter-Math leaves most of
its upright alphabets without one at all --- which is why
\verb|\hat{\mathbf{R}}| used to need a hand-tuned \verb|\skew|.

\texttt{gwbook.cls} overwrites that number for every upright glyph with the
centre of the glyph's bounding box, and leaves every slanted one exactly as the
designer set it. This sheet is the check. It is published twice from one
source, once per typeface: \mbox{\texttt{gwbook-accents.pdf}} in TeX Gyre
Schola and \mbox{\texttt{gwbook-accents-charter.pdf}} in XCharter.

\section{Upright alphabets}
\label{sec:upright}

These are the ones the class re-centres. On every letter the peak of the hat
should sit over the middle of the letter.

\medskip
\gwbcaps{\symup}{symup}
\gwblows{\symup}{}
\medskip
\gwbcaps{\symbf}{symbf}
\gwblows{\symbf}{}
\medskip
\gwbcaps{\symsf}{symsf}
\gwblows{\symsf}{}
\medskip
\gwbcaps{\symbfsf}{symbfsf}
\gwblows{\symbfsf}{}
\medskip
\gwbcaps{\symtt}{symtt}
\gwblows{\symtt}{}

\bigskip
\gwbrowstart{digits}%
\foreach \gwbL in {0,...,9}{$\hat{\gwbL}$\hskip 0.3em}\quad
\foreach \gwbL in {0,...,9}{$\hat{\symbf{\gwbL}}$\hskip 0.3em}\par

\medskip
\gwbrowstart{symup Gk}\gwbgreek{\symup}

\medskip
\gwbrowstart{symbfup Gk}\gwbgreek{\symbfup}

\section{Slanted alphabets}
\label{sec:slanted}

Untouched, for comparison: here the accent leans with the letter, which is the
font designer's judgement and not something to correct. A slanted letter's top
is to the right of its middle, and that is where the accent goes.

\medskip
\gwbcaps{\symit}{symit}
\gwblows{\symit}{}
\medskip
\gwbcaps{\symbfit}{symbfit}
\gwblows{\symbfit}{}

\medskip
\gwbrowstart{symbf Gk}\gwbgreek{\symbf}

\noindent
(\verb|\symbf| on a lowercase Greek letter gives the bold \emph{italic} shape;
that is the TeX convention, and \verb|\symbfup| is the upright one.)

\section{Reading the rule}
\label{sec:rule}

Whether an accent is centred is hard to judge by eye on a letter like
\(\symbf{L}\) or \(\symbf{J}\), whose ink is nowhere near symmetric. The
hairline below marks the centre of the base letter's ink. On the upright rows
the peak of every hat should sit on it; on the italic row it should not, and
by how much is the designer's call.

\medskip
\gwbcapsruled{\symup}{symup}
\gwblowsruled{\symup}{}
\medskip
\gwbcapsruled{\symbf}{symbf}
\gwblowsruled{\symbf}{}
\medskip
\gwbcapsruled{\symit}{symit}
\gwblowsruled{\symit}{}

\section{Every accent shape}
\label{sec:shapes}

All of them read the same attachment point, so fixing it fixes all of them at
once. Each column is one accent over an upright bold letter, a wide one, a
narrow one and a descender.

% Each accent keeps its name next to it, and each name-and-sample pair is one
% \mbox so a line break can only fall between pairs.
\newcommand{\gwbacc}[3]{%
  \mbox{{\footnotesize\ttfamily#2}\hskip 0.35em
        $#1{#3{R}}\;#1{#3{W}}\;#1{#3{l}}\;#1{#3{q}}$}%
  \hspace{1.1em plus 0.5em}}
\newcommand{\gwbaccrow}[1]{%
  \par\medskip{\raggedright\noindent
  \gwbacc{\hat}{hat}{#1}\gwbacc{\check}{check}{#1}\gwbacc{\tilde}{tilde}{#1}%
  \gwbacc{\acute}{acute}{#1}\gwbacc{\grave}{grave}{#1}\gwbacc{\dot}{dot}{#1}%
  \gwbacc{\ddot}{ddot}{#1}\gwbacc{\breve}{breve}{#1}\gwbacc{\bar}{bar}{#1}%
  \gwbacc{\vec}{vec}{#1}\gwbacc{\mathring}{mathring}{#1}\par}}

\gwbaccrow{\symbf}

\medskip
\noindent
The same over upright letters:

\gwbaccrow{\symup}

\section{Wide accents}
\label{sec:wide}

A wide accent grows to fit and is centred on the whole group, so it never
depended on the attachment point in the first place. Shown to confirm nothing
here regressed.

\[
  \widehat{AB}\quad \widehat{\symbf{ABC}}\quad \widetilde{xyz}\quad
  \overline{\symbf{R}}\quad \overrightarrow{PQ}\quad
  \overbrace{a+b+c}
\]

\section{The alphabet switches}
\label{sec:alphabets}

\verb|\mathrm|, \verb|\mathit| and \verb|\mathbf| are the Unicode alphabets
\verb|\symup|, \verb|\symit| and \verb|\symbf| in this class, so they pick up
the same fix. The two rows here must look identical.

\medskip
\noindent
\gwblabel{mathbf}%
$\hat{\mathbf{R}}\;\hat{\mathbf{A}}\;\hat{\mathbf{q}}\;\hat{\mathbf{f}}\;
 \hat{\mathbf{W}}\;\hat{\mathbf{l}}$\qquad
\gwblabel{symbf}%
$\hat{\symbf{R}}\;\hat{\symbf{A}}\;\hat{\symbf{q}}\;\hat{\symbf{f}}\;
 \hat{\symbf{W}}\;\hat{\symbf{l}}$

\medskip
\noindent
\gwblabel{mathrm}%
$\hat{\mathrm{R}}\;\hat{\mathrm{A}}\;\hat{\mathrm{q}}\;\hat{\mathrm{f}}\;
 \hat{\mathrm{W}}\;\hat{\mathrm{l}}$\qquad
\gwblabel{symup}%
$\hat{\symup{R}}\;\hat{\symup{A}}\;\hat{\symup{q}}\;\hat{\symup{f}}\;
 \hat{\symup{W}}\;\hat{\symup{l}}$

\medskip
\noindent
\verb|\mathsf| and \verb|\mathtt| stay on the document's own sans and mono
faces, which have no math data to position an accent with --- use
\verb|\symsf| and \verb|\symtt| if you need to put an accent on one:
$\hat{\mathsf{R}}$ \verb|\mathsf| versus $\hat{\symsf{R}}$ \verb|\symsf|.

\section{Vectors from the physics package}
\label{sec:physics}

\verb|physics| is loaded by default, and its unit vector \verb|\vu| is
\verb|\mathbf| applied to an accented letter, so it inherits the fix.

\[
  \vu{r}\quad \vu{n}\quad \vu{\theta}\quad \vb{E}\quad \va{v}\quad
  \grad\phi \quad \div\vb{E} \quad \curl\vb{B} \quad \laplacian\phi
\]
\[
  \dv{f}{x}\quad \pdv{f}{t}\quad \dd{x}\quad \abs{z}\quad \norm{v}\quad
  \qty(\frac{a}{b})\quad \ev{H}\quad \braket{a}{b}
\]

\noindent
\verb|physics| and \verb|unicode-math| both want \verb|\div|, \verb|\Re| and
\verb|\Im|; the class hands all three to \verb|physics| and keeps the symbols
under the names \verb|physics| gives them:
\[
  \div\vb{E} \quad \divisionsymbol \quad \Re z \quad \real z \quad
  \Im z \quad \imaginary z
\]
